\chapter{Le moment cinétique quantique}\label{ch:b3:quantum-angular-momentum}

Pointez un radiotélescope vers une région sombre du ciel et
accordez-le sur \qty{115}{GHz}~: une raie brillante et fine comme une aiguille
apparaît --- des molécules de monoxyde de carbone, à dix kelvins
au-dessus du zéro absolu, qui descendent d'un barreau d'une échelle
d'états \emph{de rotation}. Leur culbute, comme tout ce qui tourne en
mécanique quantique, est quantifiée deux fois~: la norme du moment
cinétique ne peut valoir que $\sqrt{l(l+1)}\,\hbar$, et sa composante
sur un axe choisi que $m\hbar$. Ce chapitre obtient cette double
quantification --- une fois algébriquement, à partir des \omterm{def:b3:quantum-formalism:commutator}{commutateurs}
hérités des \omterm{def:b3:hamiltonian-mechanics:poisson}{crochets de Poisson} du \cref{ch:b3:hamiltonian-mechanics},
et une fois concrètement, sous la forme des \omterm{prop:b3:quantum-angular-momentum:harmonics}{harmoniques sphériques} qui
dessinent chaque orbitale atomique. L'algèbre nous donnera plus que
nous ne demandons~: elle autorise des valeurs demi-entières qu'aucun
mouvement orbital ne peut réaliser, une case vacante que la nature
comblera deux chapitres plus loin avec le spin. En chemin nous
rencontrons les échelles de rotation des molécules --- les raies
millimétriques par lesquelles les astronomes pèsent le gaz froid des
galaxies --- et les moments magnétiques par lesquels le moment
cinétique s'est montré scindé pour la première fois.

\section{L'algèbre de la rotation}

\begin{definition}[Opérateurs de moment cinétique]\label{def:b3:quantum-angular-momentum:operators}
\index{moment cinétique!quantique}
Le moment cinétique orbital est l'opérateur $\hat{\vect L} =
\hat{\vect r}\wedge\hat{\vect p}$, de composantes $\hat L_z = \hat
x\hat p_y - \hat y\hat p_x$ et permutées. À partir de $[\hat x, \hat p_x] =
\iu\hbar$~:
\[
[\hat L_x, \hat L_y] = \iu\hbar\,\hat L_z
\quad\text{(et permutations circulaires)} , \qquad
[\hat L^2, \hat L_z] = 0 :
\]
les composantes sont mutuellement incompatibles --- aucun état n'en a
deux de bien définies --- mais le carré total est compatible avec
n'importe laquelle~: le couple $(\hat L^2, \hat L_z)$ est le choix
usuel d'étiquettes. Ces \omterm{def:b3:quantum-formalism:commutator}{commutateurs} valent exactement $\iu\hbar$ fois
les \omterm{def:b3:hamiltonian-mechanics:poisson}{crochets de Poisson} calculés dans
\cref{ex:b3:hamiltonian-mechanics:brackets}~: le dictionnaire de Dirac
à l'œuvre.
\end{definition}

\begin{theorem}[Le spectre, par la seule algèbre]\label{thm:b3:quantum-angular-momentum:spectrum}
\index{nombre quantique magnétique}
Soient $\hat J_x, \hat J_y, \hat J_z$ trois \omterm{def:b3:quantum-formalism:observable}{opérateurs hermitiens}
\emph{quelconques} obéissant aux relations de commutation ci-dessus.
Alors les \omterm{def:b3:quantum-formalism:observable}{valeurs propres} communes de $(\hat J^2, \hat J_z)$ sont
\[
\hat J^2:\ j(j+1)\hbar^2 , \qquad
\hat J_z:\ m\hbar , \quad m = -j, -j+1, \dots, +j ,
\]
où $j$ est un \emph{entier ou demi-entier} positif ou nul~: $2j + 1$
valeurs de $m$ pour chaque $j$. Les \omterm{def:b3:harmonic-oscillator:ladder}{opérateurs d'échelle} $\hat J_\pm = \hat
J_x \pm \iu\hat J_y$ font varier $m$ de $\pm1$ à $j$ fixé~:
\[
\hat J_\pm\ket{j,m} = \hbar\sqrt{j(j+1) - m(m\pm1)}\,\ket{j,m\pm1} .
\]
\end{theorem}

\begin{proof}[Démonstration partielle]
$[\hat J_z, \hat J_\pm] = \pm\hbar\hat J_\pm$~: agir avec $\hat
J_\pm$ décale la \omterm{def:b3:quantum-formalism:observable}{valeur propre} de $\hat J_z$ de $\pm\hbar$, à
$\hat J^2$ fixé (lequel commute avec tout ce qui est bâti sur les $\hat
J_i$). La norme $\|\hat J_\pm\ket{j,m}\|^2 = \hbar^2[j(j+1) -
m(m\pm1)]$ (calculée à partir de $\hat J_\mp\hat J_\pm = \hat J^2 - \hat
J_z^2 \mp \hbar\hat J_z$) doit rester positive ou nulle~: l'échelle
doit donc s'arrêter en haut à un certain $m_{\max}$ avec $m_{\max}(m_{\max}+1) =
j(j+1)$, c'est-à-dire $m_{\max} = j$, et en bas à $m_{\min} = -j$.
Monter de $-j$ à $+j$ par pas unité force $2j$ à être un entier
positif ou nul. La notation $j(j+1)\hbar^2$ pour la \omterm{def:b3:quantum-formalism:observable}{valeur propre} s'en
trouve justifiée après coup.
\end{proof}

\begin{remark}[Une case vacante au catalogue]\label{rem:b3:quantum-angular-momentum:halfinteger}
L'algèbre autorise $j = \tfrac12, \tfrac32, \dots$ --- des échelles à
nombre pair de barreaux. Le mouvement orbital, nous le montrons
ensuite, n'utilise que les entiers. La nature, cependant, ne gaspille
rien~: l'électron lui-même porte $j = \tfrac12$, sans aucune orbite
derrière lui (\cref{ch:b3:spin-two-level})~; l'algèbre obtenue ici,
inchangée, gouvernera le magnétisme atomique, les spins nucléaires et
le qubit.
\end{remark}

\section{Moment cinétique orbital~: les harmoniques sphériques}

\begin{proposition}[Harmoniques sphériques]\label{prop:b3:quantum-angular-momentum:harmonics}
\index{harmoniques sphériques}\index{nombre quantique orbital}
En coordonnées sphériques $\hat L_z = -\iu\hbar\,\partial_\varphi$, et
les fonctions propres communes de $(\hat L^2, \hat L_z)$ sur la sphère
sont les \emph{\omterm{prop:b3:quantum-angular-momentum:harmonics}{harmoniques sphériques}} $Y_l^m(\theta, \varphi)$~:
\[
\hat L^2\,Y_l^m = l(l+1)\hbar^2\,Y_l^m , \qquad
\hat L_z\,Y_l^m = m\hbar\,Y_l^m ,
\]
avec $l = 0, 1, 2, \dots$ \emph{entier} --- l'uniformité de
$\eu^{\iu m\varphi}$ force $m$ entier, donc $l$ entier. Les premières,
à la normalisation près~: $Y_0^0 = \text{cte}$ (la forme $s$, une
sphère)~; $Y_1^0 \propto \cos\theta$ et $Y_1^{\pm1}
\propto \sin\theta\,\eu^{\pm\iu\varphi}$ (les formes $p$, à deux
lobes)~; $Y_2^0 \propto 3\cos^2\theta - 1$ et ses partenaires (la
famille $d$). Parité~: $Y_l^m(-\vect r$ direction$) = (-1)^l\,Y_l^m$.
\end{proposition}
\admitted

\begin{omfigure}
\begin{tikzpicture}[font=\small, scale=0.94]
  % s orbital
  \begin{scope}
    \draw[->] (0,-1.7) -- (0,1.8);
    \node[anchor=south] at (0,1.8) {$z$};
    \draw[omDef, thick, fill=omDef!18] (0,0) circle (1.1);
    \node at (0,-2.3) {$l = 0$ ($s$)~: sphère};
  \end{scope}
  % p_z
  \begin{scope}[xshift=4.4cm]
    \draw[->] (0,-1.7) -- (0,1.8);
    \node[anchor=south] at (0,1.8) {$z$};
    \draw[omThm, thick, fill=omThm!18, domain=0:360, smooth, samples=100]
      plot ({1.45*abs(cos(\x))*sin(\x)}, {1.45*abs(cos(\x))*cos(\x)});
    \node at (0,-2.3) {$l = 1$, $m = 0$~: $|\cos\theta|$};
  \end{scope}
  % p_+-1
  \begin{scope}[xshift=8.8cm]
    \draw[->] (0,-1.7) -- (0,1.8);
    \node[anchor=south] at (0,1.8) {$z$};
    \draw[omExo, thick, fill=omExo!18, domain=0:360, smooth, samples=100]
      plot ({1.45*abs(sin(\x))*sin(\x)}, {1.45*abs(sin(\x))*cos(\x)});
    \node at (0,-2.3) {$l = 1$, $m = \pm1$~: $|\sin\theta|$};
  \end{scope}
  % d_z2
  \begin{scope}[xshift=13.2cm]
    \draw[->] (0,-1.7) -- (0,1.8);
    \node[anchor=south] at (0,1.8) {$z$};
    \draw[omProp, thick, fill=omProp!18, domain=0:360, smooth, samples=140]
      plot ({0.62*abs(3*cos(\x)*cos(\x)-1)*sin(\x)}, {0.62*abs(3*cos(\x)*cos(\x)-1)*cos(\x)});
    \node at (0,-2.3) {$l = 2$, $m = 0$};
  \end{scope}
\end{tikzpicture}

{\small Diagrammes polaires de $|Y_l^m(\theta)|$ (coupe dans un plan
contenant l'axe $z$~; la forme complète est la figure de révolution).
Ces squelettes angulaires, indépendants de tout potentiel, habilleront
chaque atome au \cref{ch:b3:hydrogen-atom}.}
\end{omfigure}

\begin{example}[Le rotateur rigide et son échelle]\label{ex:b3:quantum-angular-momentum:rotor}
Une molécule diatomique qui culbute avec le moment d'inertie $I$ a $\hat H
= \hat L^2/2I$~: les énergies
\[
E_J = \frac{\hbar^2}{2I}\,J(J+1) = B\,J(J+1) ,
\qquad J = 0, 1, 2, \dots
\]
chacune $(2J+1)$ fois dégénérée. L'absorption de photons obéit à $\Delta J =
\pm1$, si bien que les fréquences d'absorption valent $\nu_{J+1\leftarrow J} =
2B(J+1)/h$~: un peigne de raies \emph{également espacées} --- la
signature qui fait reconnaître un spectre de rotation au premier coup
d'œil. Pour le monoxyde de carbone, $B/h = \qty{57.6}{GHz}$~: la raie
fondamentale tombe à \qty{115}{GHz} ($\lambda = \qty{2.6}{mm}$), la
raie de labeur de la radioastronomie millimétrique
(\cref{pb:b3:quantum-angular-momentum:1}).
\end{example}

\section{Les potentiels centraux, complétés}

\begin{theorem}[Séparation dans un potentiel central quelconque]\label{thm:b3:quantum-angular-momentum:radial}
\index{équation radiale}\index{barrière centrifuge}
Pour $V(r)$, les états stationnaires peuvent être pris sous la forme
\[
\psi(\vect r) = \frac{u(r)}{r}\,Y_l^m(\theta, \varphi) ,
\]
où $u$ résout l'\omterm{prop:b3:schrodinger-three-dimensions:swave}{équation radiale} à une dimension
\[
-\frac{\hbar^2}{2m}\,u'' + \Big[V(r) +
\frac{\hbar^2\,l(l+1)}{2mr^2}\Big]u = E\,u , \qquad u(0) = 0 .
\]
Le problème angulaire est résolu une fois pour toutes par les $Y_l^m$~;
chaque $l$ ajoute la \emph{\omterm{thm:b3:quantum-angular-momentum:radial}{barrière centrifuge}} répulsive
$\hbar^2l(l+1)/2mr^2$ --- la version quantique du potentiel effectif
de \cref{ex:b3:lagrangian-mechanics:central} --- et chaque niveau de
$l$ donné est $(2l+1)$ fois dégénéré, car aucune force centrale ne
saurait se soucier de l'orientation de l'axe $z$. L'astuce de l'onde
$s$ de \cref{prop:b3:schrodinger-three-dimensions:swave} était la
ligne $l = 0$ de ce théorème.
\end{theorem}

\begin{proof}[Démonstration partielle]
Le laplacien se scinde en $\Delta = \tfrac1r\partial_r^2(r\,\cdot) -
\hat L^2/\hbar^2r^2$ (admis~; c'est l'énoncé selon lequel $\hat L^2$
est la partie angulaire de $-\hbar^2\Delta$). Reporter
$\psi = (u/r)Y_l^m$ et utiliser la \omterm{def:b3:quantum-formalism:observable}{valeur propre} de $\hat L^2$~:
l'équation annoncée. La dégénérescence suit de ce que $\hat H$ commute
avec les trois $\hat
L_i$, dont les \omterm{def:b3:harmonic-oscillator:ladder}{opérateurs d'échelle} déplacent $m$ sans changer l'énergie.
\end{proof}

\begin{omfigure}
\begin{tikzpicture}
  \begin{axis}[omaxis, axis x line=bottom, axis y line=left,
    width=10.6cm, height=6.0cm, xmin=0, xmax=8, ymin=-1.35, ymax=1.3,
    xlabel={$r$ (en unités de $a_0$)}, ylabel={$U_{\text{eff}}$ (en unités de $E_{\text{I}}$)},
    xlabel style={at={(axis description cs:0.5,-0.1)}, anchor=north},
    xtick={2,4,6}, ytick=\empty,
    legend style={font=\scriptsize, at={(0.97,0.95)}, anchor=north east,
      draw=none, fill=none}]
    \addplot[omDef, thick, domain=1.35:8, samples=120] {-2/x};
    \addlegendentry{$l = 0$~: Coulomb pur}
    \addplot[omThm, thick, domain=0.42:8, samples=160] {-2/x + 1/(x*x)};
    \addlegendentry{$l = 1$}
    \addplot[omExo, thick, domain=0.85:8, samples=160] {-2/x + 3/(x*x)};
    \addlegendentry{$l = 2$}
    \draw[gray, dashed] (axis cs:0,0) -- (axis cs:8,0);
  \end{axis}
\end{tikzpicture}

{\small Potentiels radiaux effectifs du problème coulombien~: la
\omterm{thm:b3:quantum-angular-momentum:radial}{barrière centrifuge} $\hbar^2l(l+1)/2mr^2$ mure l'origine pour
$l \ge 1$. Seuls les états $s$ touchent le noyau --- avec des
conséquences allant des spectres atomiques à la capture électronique
radioactive.}
\end{omfigure}

\section{Moments magnétiques~: le moment cinétique rendu visible}

\begin{proposition}[Moment magnétique orbital et effet Zeeman]\label{prop:b3:quantum-angular-momentum:magnetic}
\index{magnéton de Bohr}\index{effet Zeeman}
Une particule de charge $q$ et de masse $m$ dotée du moment cinétique
orbital $\hat{\vect L}$ porte le moment magnétique
\[
\hat{\vect\mu} = \frac{q}{2m}\,\hat{\vect L} ;
\]
pour l'électron, l'unité naturelle est le \emph{\omterm{prop:b3:quantum-angular-momentum:magnetic}{magnéton de Bohr}}
$\mu_{\text{B}} = e\hbar/2m_{\text{e}} = \qty{5.79e-5}{eV/T}$
(à comparer à \cref{exo:b3:schrodinger-three-dimensions:5}). Dans un
champ $B\vect e_z$, l'énergie $-\hat{\vect\mu}\cdot\vect B$ décale
chaque niveau de $+m\,\mu_{\text{B}}B$ (la charge de l'électron étant
négative)~: un niveau de $l$ donné se scinde en ses $2l + 1$
composantes --- l'\emph{\omterm{prop:b3:quantum-angular-momentum:magnetic}{effet Zeeman}}, première manifestation directe
de la quantification de $m$, et la raison pour laquelle $m$ s'appelle
le \omterm{thm:b3:quantum-angular-momentum:spectrum}{nombre quantique magnétique}. À $B = \qty{1}{T}$ l'écart vaut
$\qty{58}{\micro eV}$~: petit devant les énergies optiques, aisément
résolu comme un déplacement de raies spectrales --- et, lu à
l'envers, un magnétomètre~: la séparation Zeeman des raies solaires
cartographie les champs magnétiques des taches du Soleil.
\end{proposition}

\begin{proof}[Démonstration partielle]
Classiquement, une charge sur une orbite est une boucle de courant~: $\mu = IA =
(qv/2\pi r)(\pi r^2) = qL/2m$~; l'énoncé opératoriel en hérite (et
découle du couplage à $\hat{\vect A}$ de
\cref{prop:b3:lagrangian-mechanics:charge}). Le décalage en énergie
relève de la théorie des perturbations au premier ordre, anticipée ici
et justifiée au \cref{ch:b3:perturbation-theory}.
\end{proof}

\begin{omfigure}
\begin{tikzpicture}[font=\small, >=Stealth]
  \draw[->] (0,-3.0) -- (0,3.15);
  \node[anchor=south] at (0,3.15) {$z$};
  % sphere radius sqrt6
  \pgfmathsetmacro\R{2.449}
  \draw[gray, dashed] (0,0) circle (\R);
  \foreach \m/\c in {2/omDef, 1/omThm, 0/omExo, -1/omThm, -2/omDef} {
    \pgfmathsetmacro\zz{\m}
    \pgfmathsetmacro\rr{sqrt(\R*\R-\m*\m)}
    \draw[\c, thick, ->] (0,0) -- (\rr,\zz);
    \draw[gray!60, thin, dashed] (-0.05,\zz) -- (\rr,\zz);
    \node[anchor=west, \c] at ({\rr+0.1},\zz) {$m = \m$};
  }
  \node[anchor=east, font=\footnotesize] at (-0.15,2.0) {$L_z = m\hbar$};
  \node[anchor=north, font=\footnotesize] at (1.3,-2.75) {$\|\vect L\| = \sqrt{6}\,\hbar$};
\end{tikzpicture}

{\small La «~quantification spatiale~» pour $l = 2$~: le vecteur
moment cinétique a pour longueur $\sqrt{l(l+1)}\,\hbar = \sqrt6\,\hbar$
mais ne peut offrir à l'axe $z$ que les cinq projections $m\hbar$ ---
jamais sa longueur entière~: les composantes étant incompatibles,
$\vect L$ ne peut jamais se coucher exactement sur un axe.}
\end{omfigure}

\begin{method}[Le moment cinétique en pratique]\label{met:b3:quantum-angular-momentum:method}
(1) Étiqueter les états par $(l, m)$ --- ou $(j, m)$ pour l'algèbre
abstraite --- et ne jamais demander deux composantes à la fois. (2)
Éléments de matrice~: utiliser les formules d'échelle~; $\hat L_x = (\hat L_+ + \hat
L_-)/2$. (3) Potentiel central~: citer les $Y_l^m$, ne résoudre que
l'\omterm{prop:b3:schrodinger-three-dimensions:swave}{équation radiale} avec la \omterm{thm:b3:quantum-angular-momentum:radial}{barrière centrifuge}. (4) Énergies de
rotation~: $BJ(J+1)$, raies à $2B(J+1)$ --- en tirer $B$, donc les
longueurs de liaison, à partir des espacements égaux. (5) Questions
magnétiques~: convertir les moments cinétiques en moments à raison de
$\mu_{\text{B}}$ par $\hbar$, les énergies à raison de
$\mu_{\text{B}}B$ par unité de $m$.
\end{method}

\section{Exercices}

\begin{exercise}[$\star$]\label{exo:b3:quantum-angular-momentum:1}
(a) Obtenir $[\hat L_x, \hat L_y] = \iu\hbar\hat L_z$ à partir des
\omterm{def:b3:quantum-formalism:commutator}{commutateurs} canoniques. (b) Montrer que $[\hat L^2, \hat L_z] = 0$.
(c) Montrer que $[\hat L_z, \hat x] = \iu\hbar\hat y$~: qu'engendre
$\hat L_z$~? (d) Pourquoi les trois composantes n'admettent-elles
aucune base propre commune --- sauf pour un état particulier (lequel)~?
\end{exercise}

\begin{exercise}[$\star$]\label{exo:b3:quantum-angular-momentum:2}
Pour $l = 1$, dans la base $\{\ket{1,1}, \ket{1,0}, \ket{1,-1}\}$~:
(a) écrire la matrice de $\hat L_z$~; (b) utiliser la formule
d'échelle pour écrire $\hat L_+$ et $\hat L_-$~; (c) assembler
$\hat L_x$ et vérifier que ses \omterm{def:b3:quantum-formalism:observable}{valeurs propres} sont $\hbar, 0, -\hbar$~; (d) un état de $L_x =
+\hbar$ est mesuré selon $z$~: donner les probabilités des trois
issues.
\end{exercise}

\begin{exercise}[$\star$]\label{exo:b3:quantum-angular-momentum:3}
Monoxyde de carbone~: longueur de liaison \qty{0.113}{nm}, masse
réduite $\qty{1.14e-26}{kg}$. (a) Calculer $I$ et $B = \hbar^2/2I$ en
meV. (b) La fréquence et la longueur d'onde de la raie
$J = 1 \leftarrow 0$. (c) Les deux raies suivantes. (d) Une astronome
mesure l'espacement du peigne à cinq chiffres~: quelle grandeur
moléculaire obtient-elle, et avec quelle précision~?
\end{exercise}

\begin{exercise}[$\star$]\label{exo:b3:quantum-angular-momentum:4}
Une particule est dans un état de $l = 2$. (a) Énumérer les issues
possibles d'une mesure de $L_z$. (b) Le plus petit angle entre
$\vect L$ et l'axe $z$, avec $\cos\theta = m/\sqrt{l(l+1)}$~:
l'évaluer. (c) Pourquoi cet angle ne peut-il jamais être nul~? (d)
Montrer que l'angle minimal tend vers zéro quand $l \to \infty$~: les
vecteurs classiques retrouvés.
\end{exercise}

\begin{exercise}[$\star\star$]\label{exo:b3:quantum-angular-momentum:5}
La partie angulaire de $-\hbar^2\Delta$ est $\hat L^2$, avec
\[
\hat L^2 = -\hbar^2\Big[\frac{1}{\sin\theta}\,\partial_\theta
(\sin\theta\,\partial_\theta) +
\frac{1}{\sin^2\theta}\,\partial_\varphi^2\Big] .
\]
(a) Vérifier que $Y_1^0 \propto \cos\theta$ a pour \omterm{def:b3:quantum-formalism:observable}{valeur propre} de
$\hat L^2$ la valeur $2\hbar^2$. (b) Le vérifier de même pour
$Y_1^{\pm1} \propto \sin\theta\,
\eu^{\pm\iu\varphi}$, ainsi que leurs \omterm{def:b3:quantum-formalism:observable}{valeurs propres} de $\hat L_z$.
(c) Vérifier les parités. (d) Pourquoi doit-on avoir
$\braket{Y_1^0}{Y_1^1} = 0$ sans calculer la moindre intégrale~?
\end{exercise}

\begin{exercise}[$\star\star$]\label{exo:b3:quantum-angular-momentum:6}
Quelle raie de rotation est la plus brillante~? La population du
niveau $J$ est $\propto (2J+1)\,\eu^{-BJ(J+1)/k_{\text{B}}T}$. (a)
Expliquer les deux facteurs. (b) Montrer que le maximum est près de $J_{\max} \approx
\sqrt{k_{\text{B}}T/2B} - \tfrac12$. (c) Pour CO à \qty{10}{K}
(un nuage sombre) puis à \qty{300}{K}~: quelles raies dominent~? (d)
En inversant~: une échelle de CO observée dont le maximum est en
$J = 7$ trahit quelle température~?
\end{exercise}

\begin{exercise}[$\star\star$]\label{exo:b3:quantum-angular-momentum:7}
\omterm{prop:b3:quantum-angular-momentum:magnetic}{Effet Zeeman} normal. Une raie spectrale à \qty{500}{nm} provient d'une
transition $l = 2 \to l = 1$ dans un champ $B = \qty{2}{T}$~; on
ignore le spin (ce qui vaut pour des états «~singulets~» particuliers).
(a) Représenter les niveaux scindés. (b) Avec la règle de sélection
$\Delta m = 0, \pm1$, montrer que seules \emph{trois} positions de
raie apparaissent, à $0, \pm\mu_{\text{B}}B/h$. (c) Calculer l'écart en
GHz et en picomètres de longueur d'onde. (d) La plupart des raies
réelles se scindent en plus de trois composantes (Zeeman
«~anormal~»)~: de quel ingrédient manquant, porté par l'électron
lui-même, cela a-t-il été historiquement la preuve~?
\end{exercise}

\begin{exercise}[$\star\star$]\label{exo:b3:quantum-angular-momentum:8}
Le mur centrifuge. Pour l'hydrogène ($V = -e^2/4\pi\varepsilon_0r$),
écrire $U_{\text{eff}}$ pour $l = 1$ en unités de $E_{\text{I}}$ et
$a_0$. (a) Situer son minimum et sa valeur. (b) Montrer que $U_{\text{eff}} >
0$ pour $r < \dots$ --- trouver le rayon en deçà duquel la barrière
domine. (c) Comparer $|\psi|^2$ près de $r = 0$ pour les états $s$ et
$p$ ($u \sim r^{l+1}$, admis)~: quels états «~touchent~» le noyau~? (d)
La capture électronique --- un noyau avalant l'un des électrons de son
atome --- procède très majoritairement depuis les couches $s$~:
expliquer en une phrase.
\end{exercise}

\begin{exercise}[$\star\star$]\label{exo:b3:quantum-angular-momentum:9}
Sur $\ket{l, m}$~: (a) montrer que $\langle\hat L_x\rangle = \langle\hat
L_y\rangle = 0$~; (b) montrer que $\langle\hat L_x^2\rangle = \langle\hat
L_y^2\rangle = \tfrac12[l(l+1) - m^2]\hbar^2$~; (c) vérifier la
\omterm{thm:b3:quantum-formalism:uncertainty}{relation d'indétermination} $\Delta L_x\,\Delta L_y \ge \tfrac\hbar2
|\langle\hat L_z\rangle|$ sur l'état étiré $m = l$~; (d) interpréter
l'image du «~cône~» de la figure à la lumière de (a) et (b).
\end{exercise}

\begin{exercise}[$\star\star\star$]\label{exo:b3:quantum-angular-momentum:10}
Bandes de vibration-rotation. Une transition vibrationnelle
infrarouge ($\hbar\omega_0$) d'une molécule diatomique change $J$ de
$\pm1$ simultanément. (a) Montrer que les raies d'absorption tombent en $h\nu =
\hbar\omega_0 + 2B(J+1)$ (la branche $R$, de $J \to J+1$) et
$h\nu = \hbar\omega_0 - 2BJ$ (la branche $P$, $J \to J-1$), $J \ge
1$. (b) Pourquoi y a-t-il un \emph{trou} en $\hbar\omega_0$ même~? (c)
Représenter la bande~: deux peignes encadrant une raie centrale
manquante, l'intensité de chaque raie suivant
\cref{exo:b3:quantum-angular-momentum:6}. (d) La bande à
\qty{15}{\micro m} de \cref{pb:b3:harmonic-oscillator:1} a exactement
cette anatomie~: qu'est-ce qui fixe sa largeur totale, et donc les
«~ailes~» qui rendent efficace le dioxyde de carbone ajouté~?
\end{exercise}

\begin{exercise}[$\star\star\star$]\label{exo:b3:quantum-angular-momentum:11}
Achever l'algèbre. (a) À partir de $\hat J_\mp\hat J_\pm = \hat J^2 -
\hat J_z^2 \mp \hbar\hat J_z$, calculer $\|\hat J_\pm\ket{j,m}\|^2$
et les coefficients d'échelle. (b) Montrer que si l'échelle ne
s'arrêtait pas, une norme deviendrait négative~: exhiber l'état
coupable. (c) Conclure que $m_{\max} - m_{\min} = 2j$ doit être un
entier positif ou nul et énumérer les $j$ permis. (d) Où exactement
l'argument \emph{échoue}-t-il à exclure les demi-entiers --- et quelle
exigence supplémentaire (l'uniformité sur la sphère) les exclut-elle
pour le seul mouvement orbital~?
\end{exercise}

\begin{exercise}[$\star\star\star$]\label{exo:b3:quantum-angular-momentum:12}
Einstein--de Haas~: le moment cinétique est mécanique. Un cylindre de
fer (rayon $a = \qty{1.0}{cm}$, masse volumique $\rho = \qty{7.9e3}{kg/m^3}$,
$n = \qty{8.5e28}{atomes/m^3}$) est suspendu à un fil de torsion~;
renverser son aimantation retourne environ un $\hbar$ de moment
cinétique par atome. (a) Par conservation, la tige doit reculer~:
calculer le moment cinétique transféré par unité de volume. (b) Avec le moment d'inertie de la tige $\tfrac12 M
a^2$, montrer que la tige acquiert $\omega = 2n\hbar/\rho a^2$ et
l'évaluer. (c) L'expérience de 1915 entretenait le renversement à la
résonance du fil pour accumuler les minuscules coups~: estimer
l'amplitude d'oscillation après $Q \sim 100$ renversements résonnants,
en prenant votre $\omega$ par coup. (d) Le rapport mesuré
moment magnétique sur moment cinétique valait le double de la
prédiction orbitale $e/2m_{\text{e}}$~: qu'annonçait ce facteur deux~?
\end{exercise}

\section{Problème~: peser les galaxies à 2,6 millimètres}

\begin{problem}[{Problème du week-end --- monoxyde de carbone,
radioastronomie et univers froid}]\label{pb:b3:quantum-angular-momentum:1}
L'essentiel de la matière qui fabrique les étoiles d'une galaxie est
de l'hydrogène moléculaire froid --- et il est invisible~: H$_2$,
symétrique, n'a ni dipôle ni raies de rotation aux températures des
nuages froids. La solution de l'astronomie est son fidèle compagnon.
Une molécule de CO pour $\num{e4}$ H$_2$, avec son échelle à
\qty{115}{GHz}, illumine chaque nuage moléculaire du ciel. Données~:
pour CO, $B/h = \qty{57.6}{GHz}$ ($B =
\qty{0.238}{meV}$)~; longueur de liaison $r_0 = \qty{0.113}{nm}$~;
masse réduite $\mu = \qty{1.14e-26}{kg}$~; $k_{\text{B}}T$ à
\qty{10}{K} vaut \qty{0.862}{meV}.

\medskip\noindent\textbf{Partie I --- Le rotateur quantique.}
\begin{enumerate}
  \item À partir de $\hat H = \hat L^2/2I$, donner les énergies $E_J$
        et leurs dégénérescences.
  \item Vérifier $B = \hbar^2/2I$ pour CO à partir de $r_0$ et $\mu$.
  \item Avec la règle de sélection $\Delta J = \pm1$, obtenir les
        fréquences d'absorption $\nu_{J+1\leftarrow J} = 2B(J+1)/h$
        et donner les trois premières.
  \item Pourquoi les raies sont-elles \emph{également espacées} ---
        et que livre un espacement mesuré à l'astronome (via $I$)~?
  \item Calculer la longueur d'onde de la raie fondamentale, et
        expliquer pourquoi l'observer exige des sites hauts et secs
        (qu'est-ce qui absorbe les ondes millimétriques dans notre
        propre atmosphère --- rappeler la voisine de la molécule
        de \cref{pb:b3:harmonic-oscillator:1}, l'eau)~?
  \item Le photon de la raie $1 \to 0$ emporte $\hbar$ de moment
        cinétique~: pourquoi $\Delta J = \pm1$ n'est-il pas une simple
        règle empirique mais une loi de conservation~?
\end{enumerate}

\medskip\noindent\textbf{Partie II --- Pourquoi CO et pas H$_2$.}
\begin{enumerate}[resume]
  \item H$_2$ est homonucléaire~: dites pourquoi elle n'émet aucun
        rayonnement dipolaire de rotation.
  \item H$_2$ est aussi \emph{légère}~: calculer la constante de
        rotation de H$_2$ ($\mu = m_{\text{p}}/2$,
        $r_0 = \qty{0.074}{nm}$) et l'énergie de sa première
        transition accessible~; comparer à $k_{\text{B}}T$ à
        \qty{10}{K}~: même par des portes dérobées quadrupolaires,
        H$_2$ froid pourrait-il rayonner~?
  \item Le premier barreau de CO coûte $2B = \qty{0.48}{meV}$ contre
        $k_{\text{B}}T = \qty{0.86}{meV}$~: l'échelle est-elle vivante
        à \qty{10}{K}~?
  \item Calculer le rapport de populations $n_1/n_0$ à \qty{10}{K}
        (dégénérescence comprise).
  \item Près de quel $J$ la population culmine-t-elle à \qty{10}{K}
        (utiliser $J_{\max} \approx \sqrt{k_{\text{B}}T/2B} - \tfrac12$)~?
  \item Résumer en une phrase le marché conclu~: ce que fournit CO,
        ce qu'il faut supposer du rapport CO sur H$_2$.
\end{enumerate}

\medskip\noindent\textbf{Partie III --- Lire le ciel.}
\begin{enumerate}[resume]
  \item La raie d'un nuage arrive à \qty{115.156}{GHz} au lieu des
        \qty{115.271}{GHz} du laboratoire~: calculer la vitesse du
        nuage sur la ligne de visée, et son sens.
  \item La raie est \emph{élargie} par effet Doppler à cause de
        mouvements internes de $\pm\qty{2}{km/s}$~: calculer la
        largeur de raie en MHz, et comparer à la largeur thermique
        attendue à \qty{10}{K}
        ($v_{\text{th}} \sim \sqrt{k_{\text{B}}T/m_{\text{CO}}}
        \approx \qty{55}{m/s}$)~: qu'est-ce qui domine
        l'élargissement~?
  \item Cartographier le décalage Doppler de CO à travers le disque
        d'une galaxie spirale donne sa courbe de rotation $v(r)$. Ces
        courbes restent \emph{plates} bien au-delà du disque
        visible~: rappeler (volume de première année, gravitation) ce
        que $v(r)$ devrait faire autour d'une masse concentrée au
        centre, et dites ce qu'implique cette platitude.
  \item Le rapport de brillance de la raie $2\to1$ à celle de la raie
        $1\to0$ mesure les populations des niveaux~: quelle unique
        grandeur physique du nuage livre-t-il~?
  \item ALMA résout CO dans des galaxies dont la lumière est partie
        quand l'univers était jeune~; la fréquence observée de la raie
        $1 \to 0$ d'une galaxie de «~décalage vers le rouge $z = 1$~»
        arrive à la moitié de sa valeur au repos --- à quelle
        fréquence, et (rappeler
        \cref{ex:b3:relativistic-kinematics:redshift}) qu'est-ce qui
        l'a étirée~?
  \item À partir d'une luminosité de raie CO mesurée, les astronomes
        annoncent des masses de nuages en masses solaires~: énumérer
        la chaîne de conversions (photons de la raie $\to$ nombre de
        CO $\to$ masse de H$_2$) et le maillon le plus faible.
\end{enumerate}

\medskip\noindent\textbf{Partie IV --- Les pouponnières froides.}
\begin{enumerate}[resume]
  \item Les cœurs de formation stellaire sont à \qty{10}{K}~: par la
        loi de Wien, leur lueur thermique culmine vers
        \qty{290}{\micro m} --- invisible pour tout télescope
        optique. En une phrase~: pourquoi l'échelle de rotation
        est-elle le \emph{seul} thermomètre et la seule balance
        qu'offre un tel cœur~?
  \item Le niveau $J = 1$ se trouve \qty{5.5}{K} (en unités de
        température) au-dessus du fondamental~: pourquoi cela fait-il
        de la raie $1\to0$ un thermomètre exquis précisément dans le
        régime des \qty{10}{K} (plutôt que, disons, une raie optique
        à \qty{2}{eV} $\sim \qty{23000}{K}$)~?
  \item Pourquoi un cœur qui \emph{s'effondre} finit-il par ne plus
        être visible en CO ($1\to0$) --- penser à ce que la forte
        densité et la poussière font à la raie et à son évasion.
  \item Les nuages moléculaires montrent aussi des raies de NH$_3$,
        HCN et de dizaines d'autres molécules, chacune avec son
        propre $B$ et son dipôle~: que permet de démêler la richesse
        de ce «~cadran radio chimique~»~?
  \item Tout le ciel luit à \qty{2.7}{K} (le fond diffus
        cosmologique)~: qu'advient-il du contraste des raies CO d'un
        nuage à mesure que sa propre température s'approche de
        \qty{2.7}{K}, et pourquoi est-ce un plancher dur pour ce
        thermomètre~?
  \item Une antenne millimétrique doit garder sa forme à environ
        $\lambda/20$~: calculer cette tolérance à \qty{115}{GHz}, et
        comparer à celle qu'un miroir optique (\qty{500}{nm}) doit
        tenir --- pourquoi les antennes de douze mètres d'ALMA
        pouvaient-elles être bâties en plein air~?
  \item Résumer le résultat nommé~: une échelle $BJ(J+1)$ avec $B/h =
        \qty{57.6}{GHz}$, peuplée à dix kelvins et lue à
        \qty{2.6}{mm}, est la façon dont on mesure la masse, la
        température, la vitesse et la rotation de l'univers froid ---
        et l'indice de la matière noire autour des galaxies spirales.
\end{enumerate}
\end{problem}
